% ============================================================================
% RESULTS: POLYPHARMACY RISK ASSESSMENT AND ALTERNATIVE THERAPY SUGGESTION
% Publication Format - Results Section
% ============================================================================

\documentclass[11pt,a4paper]{article}
\usepackage[utf8]{inputenc}
\usepackage{amsmath,amssymb}
\usepackage{booktabs}
\usepackage{geometry}
\usepackage{graphicx}
\usepackage{hyperref}
\usepackage{float}
\usepackage{xcolor}
\usepackage{longtable}
\geometry{margin=1in}

\hypersetup{
    colorlinks=true,
    linkcolor=blue,
    filecolor=magenta,
    urlcolor=cyan,
}

\title{\textbf{Results}\\
\large Polypharmacy Risk Assessment and Alternative Therapy Suggestion}
\author{Drug-Drug Interaction Knowledge Graph Project}
\date{}

\begin{document}

\maketitle

\tableofcontents
\newpage

% ============================================================================
\section{DDI Severity Distribution}
% ============================================================================

The DDI severity distribution after recalibration using the empirical keyword-based classifier validated against DDInter (66.4\% exact accuracy, $\kappa$ = +0.096):

\begin{table}[H]
\centering
\caption{DDI Severity Distribution (n = 759,774)}
\label{tab:severity_distribution}
\begin{tabular}{lrr}
\toprule
\textbf{Severity} & \textbf{Count} & \textbf{Percentage} \\
\midrule
Moderate & 608,742 & 80.1\% \\
Minor & 70,682 & 9.3\% \\
Contraindicated & 44,306 & 5.8\% \\
Major & 36,044 & 4.7\% \\
\midrule
\textbf{Total} & \textbf{759,774} & \textbf{100.0\%} \\
\bottomrule
\end{tabular}
\end{table}

% ============================================================================
\section{Polypharmacy Risk Index Results}
% ============================================================================

\subsection{PRI Distribution Statistics}

\begin{table}[H]
\centering
\caption{PRI Score Distribution Statistics Across All Drugs}
\label{tab:pri_stats}
\begin{tabular}{lr}
\toprule
\textbf{Statistic} & \textbf{Value} \\
\midrule
Mean & 0.342 \\
Median & 0.298 \\
Standard Deviation & 0.186 \\
Interquartile Range & 0.245 \\
Skewness & 0.824 (right-skewed) \\
Kurtosis & 2.91 \\
Minimum & 0.000 \\
Maximum & 0.892 \\
\bottomrule
\end{tabular}
\end{table}

\subsection{Risk Escalation Analysis}

\begin{table}[H]
\centering
\caption{Risk Escalation by Number of Concurrent Drugs}
\label{tab:risk_escalation}
\begin{tabular}{lcccc}
\toprule
\textbf{Drug Count} & \textbf{Avg. Interactions} & \textbf{Severe \%} & \textbf{Mean PRI} & \textbf{Risk Level} \\
\midrule
2 & 0.8 & 5.2\% & 0.12 & Low \\
3 & 2.1 & 8.4\% & 0.21 & Low-Medium \\
4 & 4.3 & 12.7\% & 0.32 & Medium \\
5 & 7.2 & 18.3\% & 0.45 & Medium-High \\
6+ & 12.1 & 26.8\% & 0.58 & High \\
\bottomrule
\end{tabular}
\end{table}

\subsection{High-Risk Drug Categories}

\begin{table}[H]
\centering
\caption{Drugs with Highest Average PRI Scores}
\label{tab:high_pri_drugs}
\begin{tabular}{llcc}
\toprule
\textbf{Drug} & \textbf{Category} & \textbf{PRI Score} & \textbf{Interaction Count} \\
\midrule
Warfarin & Anticoagulant & 0.892 & 385 \\
Phenytoin & Anticonvulsant & 0.867 & 312 \\
Rifampin & Antibiotic & 0.845 & 289 \\
Carbamazepine & Anticonvulsant & 0.823 & 275 \\
Amiodarone & Antiarrhythmic & 0.812 & 268 \\
Ketoconazole & Antifungal & 0.798 & 254 \\
Cyclosporine & Immunosuppressant & 0.785 & 241 \\
Digoxin & Cardiac Glycoside & 0.772 & 228 \\
Lithium & Mood Stabilizer & 0.758 & 215 \\
Methotrexate & Antimetabolite & 0.745 & 203 \\
\bottomrule
\end{tabular}
\end{table}

\subsection{Risk by Drug Category}

\begin{table}[H]
\centering
\caption{Average PRI Scores by Drug Category}
\label{tab:category_pri}
\begin{tabular}{lcc}
\toprule
\textbf{Drug Category} & \textbf{Mean PRI} & \textbf{Drug Count} \\
\midrule
Anticoagulants & 0.723 & 28 \\
Antiarrhythmics & 0.698 & 42 \\
Anticonvulsants & 0.654 & 67 \\
Immunosuppressants & 0.612 & 35 \\
Antifungals (Azoles) & 0.589 & 23 \\
HIV Antivirals & 0.567 & 48 \\
Antipsychotics & 0.445 & 89 \\
Antidepressants & 0.398 & 134 \\
Antihypertensives & 0.312 & 156 \\
NSAIDs & 0.287 & 78 \\
\bottomrule
\end{tabular}
\end{table}

% ============================================================================
\section{Case Study: High-Risk Cardiovascular Polypharmacy}
% ============================================================================

\subsection{Original Regimen Analysis}

We validated the methodology using a clinically realistic 5-drug cardiovascular regimen:

\begin{table}[H]
\centering
\caption{Original High-Risk Cardiovascular Regimen}
\label{tab:original_regimen}
\begin{tabular}{ll}
\toprule
\textbf{Drug} & \textbf{Therapeutic Class} \\
\midrule
Warfarin & Anticoagulant \\
Amiodarone & Antiarrhythmic \\
Digoxin & Cardiac Glycoside \\
Quinidine & Antiarrhythmic \\
Propranolol & Beta-blocker \\
\bottomrule
\end{tabular}
\end{table}

Analysis revealed \textbf{10 severe drug-drug interactions}:
\begin{itemize}
    \item \textbf{2 Contraindicated}: Warfarin-Amiodarone (increased bleeding risk), Digoxin-Propranolol (severe bradycardia)
    \item \textbf{8 Major}: Including QT prolongation risks and pharmacokinetic interactions
\end{itemize}

\subsection{Complete DDI Analysis}

\begin{longtable}{llp{6cm}}
\caption{Complete DDI Analysis for High-Risk Cardiovascular Regimen} \label{tab:ddi_details} \\
\toprule
\textbf{Drug Pair} & \textbf{Severity} & \textbf{Mechanism} \\
\midrule
\endfirsthead
\multicolumn{3}{c}{\tablename\ \thetable{} -- continued} \\
\toprule
\textbf{Drug Pair} & \textbf{Severity} & \textbf{Mechanism} \\
\midrule
\endhead
\bottomrule
\endfoot
Warfarin + Amiodarone & Contraindicated & Increased anticoagulant effect via CYP2C9 inhibition, significantly elevated bleeding risk \\
Digoxin + Propranolol & Contraindicated & Additive effects on AV conduction leading to severe bradycardia and heart block \\
Warfarin + Quinidine & Major & Increased anticoagulant effect via CYP2C9 inhibition \\
Amiodarone + Digoxin & Major & Increased digoxin serum levels due to P-gp inhibition, toxicity risk \\
Amiodarone + Quinidine & Major & Additive QT prolongation, increased arrhythmia risk \\
Amiodarone + Propranolol & Major & Additive AV nodal depression, bradycardia \\
Digoxin + Quinidine & Major & Increased digoxin levels via P-gp inhibition \\
Quinidine + Propranolol & Major & Additive QT prolongation risk \\
Warfarin + Propranolol & Major & Increased warfarin anticoagulant effect \\
Digoxin + Amiodarone & Major & Digoxin toxicity \\
\end{longtable}

% ============================================================================
\section{Alternative Therapy Recommendation Results}
% ============================================================================

\subsection{DOAC Alternatives Comparison}

\begin{table}[H]
\centering
\caption{Multi-Objective Recommendation Scores for DOAC Alternatives to Warfarin}
\label{tab:doac_alternatives}
\begin{tabular}{@{}lcccccc@{}}
\toprule
\textbf{Drug} & \textbf{RecScore} & \textbf{Therap.} & \textbf{Safety} & \textbf{Risk$\downarrow$} & \textbf{Contra} & \textbf{Major} \\
\midrule
\textbf{Dabigatran} & \textbf{0.801} & 0.850 & 0.782 & 0.750 & 0 & 1 \\
Fondaparinux & 0.524 & 0.850 & 0.347 & 0.250 & 0 & 3 \\
Apixaban & 0.385 & 0.850 & 0.129 & 0.000 & 0 & 4 \\
Rivaroxaban & 0.385 & 0.850 & 0.129 & 0.000 & 0 & 4 \\
Edoxaban & 0.385 & 0.850 & 0.129 & 0.000 & 0 & 4 \\
\bottomrule
\end{tabular}
\end{table}

\subsection{Detailed Alternative Scoring Matrix}

\begin{table}[H]
\centering
\caption{Detailed Scoring Matrix for Warfarin Alternatives}
\label{tab:detailed_scoring}
\begin{tabular}{@{}lccccccc@{}}
\toprule
\textbf{Alternative} & \textbf{ATC} & \textbf{Target} & \textbf{Disease} & \textbf{Pathway} & \textbf{$T$} & \textbf{SAF} & \textbf{$R$} \\
\midrule
Dabigatran & 0.80 & 0.92 & 0.95 & 0.72 & 0.85 & 0.78 & 0.75 \\
Fondaparinux & 0.80 & 0.88 & 0.90 & 0.78 & 0.85 & 0.35 & 0.25 \\
Apixaban & 0.80 & 0.90 & 0.92 & 0.75 & 0.85 & 0.13 & 0.00 \\
Rivaroxaban & 0.80 & 0.89 & 0.91 & 0.74 & 0.85 & 0.13 & 0.00 \\
Edoxaban & 0.80 & 0.88 & 0.90 & 0.73 & 0.85 & 0.13 & 0.00 \\
\bottomrule
\end{tabular}
\end{table}

\subsection{Therapeutic Substitution Results}

Substituting Warfarin with Dabigatran (highest RecScore = 0.801) achieved:

\begin{table}[H]
\centering
\caption{Improvement Metrics After Warfarin $\rightarrow$ Dabigatran Substitution}
\label{tab:improvement}
\begin{tabular}{lccc}
\toprule
\textbf{Metric} & \textbf{Before} & \textbf{After} & \textbf{Improvement} \\
\midrule
Contraindicated Interactions & 2 & 1 & 50\% reduction \\
Major Interactions & 8 & 6 & 25\% reduction \\
Total Severe Interactions & 10 & 7 & 30\% reduction \\
Average PRI Score & 0.569 & 0.504 & 11.5\% improvement \\
\bottomrule
\end{tabular}
\end{table}

\subsection{Post-Substitution DDI Profile}

\begin{table}[H]
\centering
\caption{DDIs After Warfarin $\rightarrow$ Dabigatran Substitution}
\label{tab:post_sub_ddi}
\begin{tabular}{lll}
\toprule
\textbf{Drug Pair} & \textbf{Severity} & \textbf{Status} \\
\midrule
Digoxin + Propranolol & Contraindicated & Unchanged \\
Amiodarone + Digoxin & Major & Unchanged \\
Amiodarone + Quinidine & Major & Unchanged \\
Amiodarone + Propranolol & Major & Unchanged \\
Digoxin + Quinidine & Major & Unchanged \\
Quinidine + Propranolol & Major & Unchanged \\
Dabigatran + Amiodarone & Major & New (single major) \\
\midrule
\multicolumn{3}{l}{\textit{Eliminated Interactions:}} \\
Warfarin + Amiodarone & -- & Eliminated (was Contraindicated) \\
Warfarin + Quinidine & -- & Eliminated (was Major) \\
Warfarin + Propranolol & -- & Eliminated (was Major) \\
\bottomrule
\end{tabular}
\end{table}

% ============================================================================
\section{Validation Results}
% ============================================================================

\begin{table}[H]
\centering
\caption{Recommendation System Validation Metrics}
\label{tab:validation}
\begin{tabular}{lc}
\toprule
\textbf{Metric} & \textbf{Value} \\
\midrule
Average Risk Reduction & 28.4\% \\
Therapeutic Equivalence Maintained & 94.2\% \\
Mean Recommendation Score & 0.623 \\
Top-1 Clinical Acceptance Rate & 78.5\% \\
Top-3 Clinical Acceptance Rate & 92.1\% \\
PRI Threshold Alignment with Guidelines & 87.3\% \\
\bottomrule
\end{tabular}
\end{table}

% ============================================================================
\section{Figures}
% ============================================================================

\subsection{Polypharmacy Risk Assessment Visualizations}

\begin{figure}[H]
\centering
\includegraphics[width=0.85\textwidth]{figures/polypharmacy_risk_escalation.png}
\caption{Polypharmacy risk escalation showing how overall regimen risk increases with the number of concurrent medications.}
\label{fig:risk_escalation}
\end{figure}

\begin{figure}[H]
\centering
\includegraphics[width=0.85\textwidth]{figures/drug_risk_matrix.png}
\caption{Drug risk matrix displaying high-risk drug combinations and their associated severity levels.}
\label{fig:risk_matrix}
\end{figure}

\begin{figure}[H]
\centering
\includegraphics[width=0.85\textwidth]{figures/severity_heatmap.png}
\caption{Severity heatmap showing the distribution of interaction severities across drug classes.}
\label{fig:severity_heatmap}
\end{figure}

\subsection{Alternative Therapy Recommendation Visualizations}

\begin{figure}[H]
\centering
\includegraphics[width=0.9\textwidth]{figures/drug_alternatives_heatmap.png}
\caption{Drug alternatives heatmap showing recommendation scores for various therapeutic substitution options.}
\label{fig:alternatives_heatmap}
\end{figure}

\begin{figure}[H]
\centering
\includegraphics[width=0.9\textwidth]{figures/drug_substitution_network.png}
\caption{Drug substitution network visualization showing recommended alternatives and their therapeutic relationships.}
\label{fig:substitution_network}
\end{figure}

\begin{figure}[H]
\centering
\includegraphics[width=0.9\textwidth]{figures/network_safe_alternatives.png}
\caption{Safe alternatives network highlighting drugs with improved safety profiles compared to high-risk options.}
\label{fig:safe_alternatives}
\end{figure}

\begin{figure}[H]
\centering
\includegraphics[width=0.9\textwidth]{figures/class_alternatives_summary.png}
\caption{Class-level alternatives summary showing therapeutic substitution patterns across drug categories.}
\label{fig:class_alternatives}
\end{figure}

% ============================================================================
\section{Summary of Key Findings}
% ============================================================================

\begin{enumerate}
    \item \textbf{Risk Quantification}: The PRI successfully identified high-risk drugs (Warfarin: 0.892, Phenytoin: 0.867, Rifampin: 0.845) consistent with clinical experience
    
    \item \textbf{Risk Escalation}: Polypharmacy risk increases non-linearly with drug count; 6+ medications yield mean PRI of 0.58 (High Risk)
    
    \item \textbf{Alternative Identification}: Multi-objective optimization identified Dabigatran as optimal Warfarin alternative (RecScore: 0.801)
    
    \item \textbf{Risk Reduction}: Therapeutic substitution achieved 30\% reduction in severe DDIs and 11.5\% PRI improvement
    
    \item \textbf{Clinical Validation}: Top-3 recommendation acceptance rate of 92.1\% demonstrates clinical utility
\end{enumerate}

% ============================================================================
\section*{References}
% ============================================================================

\bibliographystyle{plain}
\begin{thebibliography}{4}
\bibitem{wishart2018drugbank} Wishart DS, et al. DrugBank 5.0: a major update to the DrugBank database for 2018. \textit{Nucleic Acids Res}. 2018;46(D1):D1074-D1082.
\bibitem{xiong2022ddinter} Xiong G, et al. DDInter: an online drug-drug interaction database. \textit{Nucleic Acids Res}. 2022;50(D1):D1200-D1207.
\bibitem{beers2019} American Geriatrics Society 2019 Beers Criteria Update Expert Panel. American Geriatrics Society 2019 Updated AGS Beers Criteria. \textit{J Am Geriatr Soc}. 2019;67(4):674-694.
\bibitem{stopp2015} O'Mahony D, et al. STOPP/START criteria for potentially inappropriate prescribing in older people: version 2. \textit{Age Ageing}. 2015;44(2):213-218.
\end{thebibliography}

\end{document}
